\documentclass[10pt,a4paper]{amsart}
\usepackage[margin=19mm]{geometry}
\usepackage{amsmath,amssymb,mathtools}
\usepackage[scr=rsfs,cal=euler]{mathalfa}
\usepackage{xfrac}
\usepackage[unicode,hidelinks,bookmarks=false]{hyperref}
\newcommand{\ie}{i.e.\ }
\newcommand{\cf}{cf.\ }
\newcommand{\resp}{respectively\ }
\newcommand{\Subsection}{\subsection}
\providecommand{\nobreakdash}{\nobreak\hbox{-}}
\setlength{\emergencystretch}{2em}
\multlinegap0pt
\usepackage{bm}
\usepackage{bbm}
\usepackage{dsfont}
\usepackage{xspace}
\usepackage{cancel}
\usepackage{mathtools}
\usepackage{amsthm}
\makeatletter
\@ifundefined{theorem}{\newtheorem{theorem}{Theorem}}{}
\makeatother
\newcommand{\R}{\mathbb{R}}
\newcommand{\X}{\textbf{X}}
\newcommand{\norm}[1]{\left\lVert#1\right\rVert}
\newcommand{\tr}{\text{tr}}
\renewcommand{\P}{P}
\title[Hyper-V uniform ergodicity of Markov chains]{Hyper-V uniform ergodicity of Markov chains}
\author{Austin Brown, Kshitij Khare}
\date{Source: arXiv:2608.12738v1 (CC BY 4.0)}
\begin{document}
\maketitle

\noindent\emph{Source and licence.} Hyper-V uniform ergodicity of Markov chains. Version arXiv:2608.12738v1 (CC BY 4.0), distributed under the Creative Commons Attribution 4.0 International licence, as recorded in the arXiv licence metadata for this exact version and shown on the abstract page https://arxiv.org/abs/2608.12738v1. Full rights notice: licenses/arxiv-2608.12738-CC-BY-4.0.txt.\par\smallskip

\noindent\textit{Editorial scope.} Counted unit: the Theorem whose source label is \texttt{thm:convergence\_sharp} in the article (source0.tex lines 370--393, source label \texttt{thm:convergence\_sharp}), delivered together with the article's own complete original proof of it (source0.tex lines 395--593, one \texttt{proof} environment of 199 lines). No mathematics is written, repaired or re-derived: statement and proof are the article's own text line for line. Required context retained uncounted: eq:drift (lines 186--189); eq:local\_minorization (lines 193--196). Every definition, display and label of those lines is retained unchanged. Omitted from this package and not redistributed: 1--185, 190--192, 197--369, 394--394, 594--1421. Nothing inside a retained excerpt is omitted or altered: every retained body is delivered exactly as the article's lines are written, so no omission receipt of any kind accompanies this package. Citations: the retained excerpts carry no citation command at all, so no bibliography entry of the article is transcribed and no reference is redistributed. Reference adaptation: none was needed and none was made; every reference inside the delivered unit resolves inside it, so no unresolved reference or citation is produced. Printed slips kept literal: no printed word, symbol, hypothesis or number of the counted statement or of its proof was altered. Layout and class: the article's own class is \texttt{imsart} with packages including natbib, hyperref, graphicx, subcaption, comment; the wrapper is the lane's pinned \texttt{amsart} editorial wrapper at 10pt on A4, which restates the theorem-like environments the retained text needs and adds the article's own macro definitions that the retained text needs (\texttt{\textbackslash P}, \texttt{\textbackslash R}, \texttt{\textbackslash X}, \texttt{\textbackslash norm}, \texttt{\textbackslash tr}), with extra wrapper packages (bm, bbm, dsfont, xspace, cancel, mathtools). The delivered numbering is the unit's own. Assets: no retained span contains a figure environment, a graphics-inclusion command, a PDF-page inclusion, a table or a TikZ picture; no image file is copied into this package, the package declares no asset dependency and no image file is redistributed with it. Licence: version arXiv:2608.12738v1 of this article is distributed under the Creative Commons Attribution 4.0 International licence, as recorded in the arXiv licence metadata for this exact version and shown on the abstract page https://arxiv.org/abs/2608.12738v1. A full rights notice is delivered at licenses/arxiv-2608.12738-CC-BY-4.0.txt. Characters: every retained excerpt is pure ASCII, so the delivered engine, XeLaTeX with its default Latin Modern text font, reports no missing character.
\begin{align}
P V(x) \le K, \qquad \forall x \in \X.
\label{eq:drift}
\end{align}

\begin{align}
\inf_{\{x \in \X : V(x) \le rK\}} P(x, \cdot) \;\ge\; \alpha_r \, \nu(\cdot).
\label{eq:local_minorization}
\end{align}

\begin{theorem}[Optimal hyper-$V$ uniform ergodicity]  \label{thm:convergence_sharp}
Assume the drift condition \eqref{eq:drift} holds with a Borel measurable function $V : \X \to [0, \infty)$ and constant $K \in (0,\infty)$ and the minorization condition \eqref{eq:local_minorization} holds with a constant $\alpha_r \in (0, 1]$ and Borel probability measure $\nu$ on $\X$.
Let $\gamma_r = 1 - 1/r$, and define
\[
\lambda_{\pm}(\gamma_r, \alpha_r)
= \frac{1 - \gamma_r \wedge \alpha_r \pm \sqrt{ (1 - \gamma_r \wedge \alpha_r)^2 + 4 ( \gamma_r \wedge \alpha_r ) (1 - \alpha_r \vee \gamma_r) } }{2}.
\]
Let $c : [0, \infty) \to [0, \infty)$ be any nondecreasing concave function.
Then for all $x \in \X$ and for all integers $t \ge 0$,
\begin{align*}
&\left\| P^{t}(x, \cdot) - \Pi(\cdot) \right\|_{\mathrm{TV}}
\le \rho_t(\gamma_r, \alpha_r)
\\
&\left\| P^{t + 1}(x, \cdot) - \Pi(\cdot) \right\|_{c(V)}
\le 2 c(K) \rho_{t}(\gamma_r, \alpha_r)
\end{align*}
where
\begin{align}
\rho_{t}(\gamma_r, \alpha_r)
= \frac{(1 - \lambda_{-}(\gamma_r, \alpha_r)) \lambda_{+}(\gamma_r, \alpha_r)^t - (1 - \lambda_{+}(\gamma_r, \alpha_r)) \lambda_{-}(\gamma_r, \alpha_r)^t }{ \lambda_{+}(\gamma_r, \alpha_r) - \lambda_{-}(\gamma_r, \alpha_r) }
\label{eq:sharp_rate}
\end{align}
and $\lim_{t \to \infty} \rho_{t}(\gamma_r, \alpha_r)^{1/t} = \lambda_{+}(\gamma_r, \alpha_r)$.
\end{theorem}

\begin{proof}
It suffices to complete the proof assuming $\alpha_r \in (0, 1)$.
For positive integers $t$ and Borel sets $A \subseteq \X$, define $d_t(A) = \sup_{x, y \in A} \norm{\P^t(x, \cdot) - \P^t(y, \cdot)}_{\text{TV}}$.
Define the set $C_r = \{ x \in \X : V(x) \le  r K \}$. 
By Markov's inequality,
\[
P(x, C_r^c)
\le \frac{PV(x)}{r K}
\le \frac{K}{r K}
= 1- \gamma_r,
\qquad \forall x \in \X.
\]
This implies that 
\[
P(x, C_r) \ge \gamma_r.
\]
Let $x \in \X$ and write $\mu_x = \P(x, \cdot)$.
Then we can define the conditional measure $\mu_{x, C_r} = \P(x, \cdot \cap C_r) / \P(x, C_r)$.
For every Borel $B \subseteq \X$,
\[
\mu_x(B) - \gamma \mu_{x, C_r}(B)
= \mu_{x, C_r}(B) \mu_x(C_r) + \mu_x(B \cap C_r^c) - \gamma \mu_{x, C_r}
\ge 0.
\]
The residual probability measure is well defined where
\[
R_x = \frac{\mu_x - \gamma \mu_{x, C_r}}{1 - \gamma_r}
\]
and $\P(x, \cdot) = \gamma_r \mu_{x, C_r} + (1 - \gamma_r) R_x$.
For $x, y \in \X$, we have the simple identity
\[
\P^{t + 1}( x, \cdot ) - \P^{t + 1}( y, \cdot )
= \mu_{x} \P^{t} - \mu_{y} \P^{t}.
\]
Using the residual measure, we then have
\begin{align*}
\P^{t + 1}( x, \cdot ) - \P^{t + 1}( y, \cdot )
= \gamma_r ( \mu_{x, C_r} \P^{t} - \mu_{y, C_r} \P^{t} ) + ( 1 - \gamma_r ) ( R_x \P^{t} - R_y \P^{t} ).
\end{align*}
Taking the supremum over $x, y$, we conclude that
\[
d_{t + 1}(\X)
\le \gamma_r d_{t}(C_r) + ( 1 - \gamma_r ) d_{t}(\X).
\]
Using the minorization condition \eqref{eq:local_minorization}, for $x \in C_r$, we can then define a residual probability measure
\[
S_x(\cdot) = \frac{ \P(x, \cdot) - \alpha_r \nu(\cdot) }{ 1 - \alpha_r }.
\]
Then for all $x, y \in C_r$,
\begin{align*}
\left| \P^{t + 1}( x, \cdot ) - \P^{t + 1}( y, \cdot ) \right|
&\le ( 1 - \alpha_r ) \left| S_x \P^{t} - S_y \P^{t} \right|.
\end{align*}
Using the drift condition \eqref{eq:drift},
\[
S_x(C_r^c)
\le \frac{\P V(x)}{ r K  (1 - \alpha_r)}
\le \frac{1}{ r  (1 - \alpha_r)}
\le \frac{1 - \gamma_r}{1 - \alpha_r}.
\]
So 
\[
S_x(C_r)
\ge \frac{(\gamma_r - \alpha_r)_+}{1 - \alpha_r}.
\]
Let $\beta_r = (\gamma_r - \alpha_r)_+$ denoting the positive part.
Applying the previous argument used with $\mu_x$, then for all $x, y \in C_r$,
\begin{align*}
( 1 - \alpha_r ) \left| S_x \P^{t} - S_y \P^{t} \right|
\le ( 1 - \alpha_r ) \left( \frac{\beta_r}{1 - \alpha_r} d_{t}(C_r) + \left( 1 - \frac{\beta_r}{1 - \alpha_r} \right) d_{t}(\X) \right).
\end{align*}
Therefore,
\begin{align*}
d_{t+1}(C_r) &\le \beta_r d_{t}(C_r) + \left( 1 - \alpha_r - \beta_r \right) d_{t}(\X).
\end{align*}
Define the matrix 
\[
M_r = \begin{pmatrix}
1 - \gamma_r & \gamma_r \\
1 - \alpha_r - \beta_r  & \beta_r
\end{pmatrix}.
\]
Then with the inequality understood coordinate-wise:
\begin{align*}
\begin{pmatrix}
d_{t+1}(\X) \\
d_{t+1}(C_r)
\end{pmatrix}
&\le M_r \begin{pmatrix}
d_{t}(\X) \\
d_{t}(C_r)
\end{pmatrix}
\end{align*}
Since $M_r$ has nonnegaitve entries, then we apply this recursively to get
\begin{align*}
\begin{pmatrix}
d_{t}(\X) \\
d_{t}(C_r)
\end{pmatrix}
\le M_r^{t} \begin{pmatrix}
d_{0}(\X) \\
d_{0}(C_r)
\end{pmatrix}.
\end{align*}
Define $\rho_{r, t} = (1, 0) M_r^t (1, 1)^T$.
Then using matrix multiplication $d_{t}(\X) \le \rho_{r, t}$.
We have then shown that for every $x, y \in \X$,
\[
\norm{\P^t(x, \cdot) - \P^t(y, \cdot)}_{\text{TV}}
\le \rho_{r, t}.
\]
Therefore, integrating with respect to $\Pi$ concludes
\[
\sup_{x \in \X} \norm{\P^t(x, \cdot) - \Pi(\cdot)}_{\text{TV}}
\le \rho_{r, t}.
\]

We now extend this to uniform convergence in the weighted $V$-norm. 
Define the oscillation of a bounded measurable function $f : \X \to \R$ by
\[
\text{osc}(f) := \sup_{x,y \in \X} |f(x) - f(y)|.
\]
Now let $\psi : \X \to \R$ with $|\psi| \le c(V)$. 
By Jensen's inequality, for $x \in \X$
\[
\P |\psi|(x) 
\le \int_{\X} c(V(y)) \P(x, dy)
\le c\left( \int_{\X} V(y) \P(x, dy) \right)
\le c( K ).
\]
Then $| P\psi |$ is bounded by the drift condition, and
\[
\sup_{x \in \X} |P\psi(x)| \le c( K ).
\]
Thus,
\begin{align*}
\text{osc}(P^{t+1}\psi)
\le \rho_{r, t}  2 \sup_{x} |P\psi(x)|
\le \rho_{r, t}  2 c(K).
\end{align*}
Therefore, integrating with respect to $\Pi$
\[
\sup_{x \in \X} \norm{\P^{t + 1}(x, \cdot) - \Pi(\cdot)}_{c(V)}
\le \rho_{r, t} 2 c(K).
\]

It remains to show an explicit form of $\rho_{r, t}$ that is represented by formula \eqref{eq:sharp_rate}.
First assume $\gamma_r > \alpha_r$.
The characteristic polynomial of $M_r$ is defined by
\[
p(\lambda)
= \lambda^2 - \tr(M_r) \lambda + \det(M_r)
=  \lambda^2 - (1 - \alpha_r) \lambda - \alpha_r (1 - \gamma_r).
\]
Next, assume $\gamma_r \le \alpha_r$.
The characteristic polynomial in this case is
\[
p(\lambda)
= \lambda^2 - \tr(M_r) \lambda + \det(M_r)
= \lambda^2 - (1 - \gamma_r) \lambda - \gamma_r (1 - \alpha_r).
\]
Therefore, it suffices to complete the proof assuming $\gamma_r \le \alpha_r$. 
In this case, the characteristic polynomial has roots
\[
\lambda_{\pm} = \frac{1 - \gamma_r \pm \sqrt{ (1 - \gamma_r)^2 + 4 \gamma_r (1 - \alpha_r) } }{2}.
\]
Using the Cayley-Hamilton theorem
\[
M_R^2 = (1 - \gamma_r) M_r + \gamma_r (1 - \alpha_r) I
\]
where $I$ is the identity matrix.
So then multiplying by the row vector $(1, 0)$,
\begin{align*}
&M_r^{t + 2} = (1 - \gamma_r) M_r^{t + 1} + \gamma_r (1 - \alpha_r) M_r^t
\\
&\rho_{r, t + 2} = (1 - \gamma_r) \rho_{r, t + 1} + \gamma_r (1 - \alpha_r) \rho_{r, t}.
\end{align*}
Now, similarly both roots satisfy this relation and for $a \in [0, \infty)$
\[
a \lambda_+^t + (1 - a) \lambda_-^t
\]
satisfies the relation as well.
We can solve directly that $\rho_{r, 0} = \rho_{r, 1} = 1$.
We must solve $a$ to satisfy the initial conditions and then $\rho_{r, t} = a \lambda_+^t + (1 - a) \lambda_-^t$ for all integers $t \ge 0$.
Since $a + 1 - a = 1$, this is true if
$
a \lambda_+ + (1 - a) \lambda_- = 1.
$
This is solved by 
\[
a = \frac{1 - \lambda_-}{ \lambda_+ - \lambda_- }.
\]
Therefore,
\[
\rho_{r, t}
= \frac{(1 - \lambda_-) \lambda_+^t + (\lambda_+ - 1) \lambda_-^t }{ \lambda_+ - \lambda_- }.
\]
Since $\lambda_+ > -\lambda_-$ using Vieta's formula, $(\lambda_- / \lambda_+)^t \to 0$ as $t \to \infty$, then $\lim_{t \to \infty} \rho_{r, t}^{1/t} = \lambda_+$.
\end{proof}

\end{document}
